% GENERATED FILE. Edit canonical lessons, then run npm run generate:books.
% canonical-source-hash: fnv1a64:3b7eb759b5133447
% canonical-lessons: LA-C356-appono, LA-C356-ministro, LA-C356-triclinium, LA-C356-discumbo, LA-C356-ferculum, LA-R356-review-355-356, LA-R356-review-final-356

\chapter{Cēna appōnitur}
\label{ch:la-cena-apponitur}
\hlchaptermodality{356}

\begin{chapteropening}
\textbf{By the end of this chapter:} \emph{I can tell how a dinner was served and eaten: to serve up, to wait at table, the dining room, to recline and a course.}

Complete the last lesson of chapter 356: I can tell how a dinner was served and eaten: to serve up, to wait at table, the dining room, to recline and a course.
\end{chapteropening}

\section[appōnō, appōnere]{appōnō, appōnere — to serve up, to put before}

\label{lesson:LA-C356-appono}



\begin{quote}
Before the new one: say the Latin for to pour in, then the Latin for to strain.
\end{quote}

\subsection*{You'll want to know: appōnō, appōnere}
\textbf{appōnō, appōnere} — \textquotedblleft{}to serve up, to put before\textquotedblright{}.

\begin{quote}
Servī cēnam amīcīs appōnēbant.
\end{quote}

\begin{quote}
\emph{The slaves were serving dinner to the friends.}
\end{quote}

\emph{appōnō, appōnere, apposuī, appositum}: \emph{ad} + \emph{pōnō} (taught), to set food before someone.

\begin{cousinweb}
\textbf{appōnō} is built on \emph{pōnō}, \textquotedblleft{}to put\textquotedblright{}, with \emph{ad-}, \textquotedblleft{}to, beside\textquotedblright{} (Lewis \& Short): to put before someone.
\end{cousinweb}

\begin{grammarlens}[title={What you've built}]
One more word for telling, explaining and describing at B1.
\end{grammarlens}

\subsection*{Your turn}
\begin{itemize}
\raggedright
  \item \emph{Say it:} \emph{appōnō, appōnere}
  \item \emph{Say it:} \emph{appōnō, appōnere}, once more
  \item \emph{Say it:} a sentence of your own with \emph{appōnō, appōnere}
  \item \emph{Recall:} say the Latin for to pour in, then the Latin for to strain, then say \emph{appōnō, appōnere} again
\end{itemize}

\begin{tcolorbox}[breakable,colback=teal!4,colframe=teal!35!black,title={Before you move on}]
What does appōnō, appōnere mean? (\textquotedblleft{}To serve up, to put before\textquotedblright{}.) Say it once more.
\end{tcolorbox}
\section[ministrō, ministrāre]{ministrō, ministrāre — to wait at table, to serve}

\label{lesson:LA-C356-ministro}



\begin{quote}
Before the new one: say the Latin for to strain, then the Latin for to serve up.
\end{quote}

\subsection*{You'll want to know: ministrō, ministrāre}
\textbf{ministrō, ministrāre} — \textquotedblleft{}to wait at table, to serve\textquotedblright{}.

\begin{quote}
Puerī vīnum ministrābant.
\end{quote}

\begin{quote}
\emph{The slave boys were serving the wine.}
\end{quote}

\emph{ministrō, ministrāre, ministrāvī, ministrātum}, from \emph{minister}, 'servant'. English \emph{minister}.

\begin{cousinweb}
\textbf{ministrō} is built on \emph{minister}, \textquotedblleft{}servant, attendant\textquotedblright{} (EtymDB), itself from \emph{minus}, \textquotedblleft{}less\textquotedblright{}, as \emph{magister} is from \emph{magis}. English \emph{minister} comes from it.
\end{cousinweb}

\begin{grammarlens}[title={What you've built}]
One more word for telling, explaining and describing at B1.
\end{grammarlens}

\subsection*{Your turn}
\begin{itemize}
\raggedright
  \item \emph{Say it:} \emph{ministrō, ministrāre}
  \item \emph{Say it:} \emph{ministrō, ministrāre}, once more
  \item \emph{Say it:} a sentence of your own with \emph{ministrō, ministrāre}
  \item \emph{Recall:} say the Latin for to strain, then the Latin for to serve up, then say \emph{ministrō, ministrāre} again
\end{itemize}

\begin{tcolorbox}[breakable,colback=teal!4,colframe=teal!35!black,title={Before you move on}]
What does ministrō, ministrāre mean? (\textquotedblleft{}To wait at table, to serve\textquotedblright{}.) Say it once more.
\end{tcolorbox}
\section[triclīnium, triclīniī]{triclīnium, triclīniī — a dining room; a dining couch}

\label{lesson:LA-C356-triclinium}



\begin{quote}
Before the new one: say the Latin for to serve up, then the Latin for to wait at table.
\end{quote}

\subsection*{You'll want to know: triclīnium, triclīniī}
\textbf{triclīnium, triclīniī} — \textquotedblleft{}a dining room; a dining couch\textquotedblright{}.

\begin{quote}
Hospitēs in triclīnium intrāvērunt.
\end{quote}

\begin{quote}
\emph{The guests went into the dining room.}
\end{quote}

Neuter, a Greek loan: three couches round a table, three guests on each.

\begin{cousinweb}
\textbf{triclīnium} was borrowed from Greek \emph{triklinion}, from \emph{treis}, \textquotedblleft{}three\textquotedblright{}, and \emph{klinē}, \textquotedblleft{}couch\textquotedblright{} (EtymDB): the room with three couches.
\end{cousinweb}

\begin{grammarlens}[title={What you've built}]
One more word for telling, explaining and describing at B1.
\end{grammarlens}

\subsection*{Your turn}
\begin{itemize}
\raggedright
  \item \emph{Say it:} \emph{triclīnium, triclīniī}
  \item \emph{Say it:} \emph{triclīnium, triclīniī}, once more
  \item \emph{Say it:} a sentence of your own with \emph{triclīnium, triclīniī}
  \item \emph{Recall:} say the Latin for to serve up, then the Latin for to wait at table, then say \emph{triclīnium, triclīniī} again
\end{itemize}

\begin{tcolorbox}[breakable,colback=teal!4,colframe=teal!35!black,title={Before you move on}]
What does triclīnium, triclīniī mean? (\textquotedblleft{}A dining room; a dining couch\textquotedblright{}.) Say it once more.
\end{tcolorbox}
\section[discumbō, discumbere]{discumbō, discumbere — to recline at table}

\label{lesson:LA-C356-discumbo}



\begin{quote}
Before the new one: say the Latin for to wait at table, then the Latin for a dining room.
\end{quote}

\subsection*{You'll want to know: discumbō, discumbere}
\textbf{discumbō, discumbere} — \textquotedblleft{}to recline at table\textquotedblright{}.

\begin{quote}
Omnēs discubuērunt, et cēna incēpit.
\end{quote}

\begin{quote}
\emph{Everyone took their places at table, and the dinner began.}
\end{quote}

\emph{discumbō, discumbere, discubuī, discubitum}. Romans ate lying on the left elbow.

\begin{cousinweb}
\textbf{discumbō} is built on \emph{-cumbō}, \textquotedblleft{}to lie down\textquotedblright{}, with \emph{dis-}, \textquotedblleft{}apart\textquotedblright{} (EtymDB): each guest lying down in his place.
\end{cousinweb}

\begin{grammarlens}[title={What you've built}]
One more word for telling, explaining and describing at B1.
\end{grammarlens}

\subsection*{Your turn}
\begin{itemize}
\raggedright
  \item \emph{Say it:} \emph{discumbō, discumbere}
  \item \emph{Say it:} \emph{discumbō, discumbere}, once more
  \item \emph{Say it:} a sentence of your own with \emph{discumbō, discumbere}
  \item \emph{Recall:} say the Latin for to wait at table, then the Latin for a dining room, then say \emph{discumbō, discumbere} again
\end{itemize}

\begin{tcolorbox}[breakable,colback=teal!4,colframe=teal!35!black,title={Before you move on}]
What does discumbō, discumbere mean? (\textquotedblleft{}To recline at table\textquotedblright{}.) Say it once more.
\end{tcolorbox}
\section[ferculum, ferculī]{ferculum, ferculī — a course, a dish; a tray}

\label{lesson:LA-C356-ferculum}



\begin{quote}
Before the new one: say the Latin for a dining room, then the Latin for to recline at table.
\end{quote}

\subsection*{You'll want to know: ferculum, ferculī}
\textbf{ferculum, ferculī} — \textquotedblleft{}a course, a dish; a tray\textquotedblright{}.

\begin{quote}
Prīmum ferculum ōva et olīvae erant.
\end{quote}

\begin{quote}
\emph{The first course was eggs and olives.}
\end{quote}

Neuter, from \emph{ferō}: a tray that carries food, and so a course. In a triumph, a \emph{ferculum} carried the spoils.

\begin{cousinweb}
\textbf{ferculum} is built on \emph{ferō}, \textquotedblleft{}to carry\textquotedblright{}, with the tool ending \emph{-culum} (EtymDB): a carrier, a tray.
\end{cousinweb}

\begin{grammarlens}[title={What you've built}]
One more word for telling, explaining and describing at B1.
\end{grammarlens}

\subsection*{Your turn}
\begin{itemize}
\raggedright
  \item \emph{Say it:} \emph{ferculum, ferculī}
  \item \emph{Say it:} \emph{ferculum, ferculī}, once more
  \item \emph{Say it:} a sentence of your own with \emph{ferculum, ferculī}
  \item \emph{Recall:} say the Latin for a dining room, then the Latin for to recline at table, then say \emph{ferculum, ferculī} again
\end{itemize}

\begin{tcolorbox}[breakable,colback=teal!4,colframe=teal!35!black,title={Before you move on}]
What does ferculum, ferculī mean? (\textquotedblleft{}A course, a dish; a tray\textquotedblright{}.) Say it once more.
\end{tcolorbox}
\section[(dialogue)]{The last stories — a review}

\label{lesson:LA-R356-review-355-356}



\begin{quote}
Say the Latin for to season and a seasoning.
\end{quote}

\subsection*{Your turn}
Say these aloud:

\begin{itemize}
\raggedright
  \item to season, a seasoning, to sprinkle, to pour in, to strain
  \item to serve up, to wait at table, a dining room, to recline at table, a course
\end{itemize}

\begin{tcolorbox}[breakable,colback=teal!4,colframe=teal!35!black,title={Before you move on}]
To season? (\textbf{condiō, condīre}.) A seasoning? (\textbf{condīmentum, condīmentī}.) To sprinkle? (\textbf{aspergō, aspergere}.) To pour in? (\textbf{īnfundō, īnfundere}.) To strain? (\textbf{cōlō, cōlāre}.) To serve up? (\textbf{appōnō, appōnere}.) To wait at table? (\textbf{ministrō, ministrāre}.) A dining room? (\textbf{triclīnium, triclīniī}.) To recline at table? (\textbf{discumbō, discumbere}.) A course? (\textbf{ferculum, ferculī}.)
\end{tcolorbox}
\section[(dialogue)]{Cēna appōnitur — once more}

\label{lesson:LA-R356-review-final-356}



\begin{quote}
Say the Latin for to serve up and to wait at table.
\end{quote}

\subsection*{Your turn}
\begin{itemize}
\raggedright
  \item \emph{Say it:} to serve up, to wait at table, a dining room, to recline at table, a course
\end{itemize}

\begin{tcolorbox}[breakable,colback=teal!4,colframe=teal!35!black,title={Before you move on}]
To serve up? (\textbf{appōnō, appōnere}.) To wait at table? (\textbf{ministrō, ministrāre}.) A dining room? (\textbf{triclīnium, triclīniī}.) To recline at table? (\textbf{discumbō, discumbere}.) A course? (\textbf{ferculum, ferculī}.)
\end{tcolorbox}
